\subsection{What is intelligence? What is AI?}
\label{sec:what-is-intelligence}
\slidesrc{Thinking machines: 4--12}

To build \emph{artificial} intelligence we should first agree on what intelligence is.
There is no single accepted answer; the slides show several complementary views.

\subsubsection{The word and a classic definition}
\slidesrc{Thinking machines: 4--6}

The word comes from Latin \emph{intellig\u{e}re}: ``to understand'', literally
``to read inside'' (\emph{intus legere}) or ``to read between'' (\emph{inter legere}),
that is, to read beyond the surface of things.

\begin{nbdef}{Intelligence (Mainstream Science on Intelligence, 1994)}
Intelligence is a very general mental capability that, among other things, involves the
ability to \textbf{reason}, \textbf{plan}, \textbf{solve problems}, \textbf{think abstractly},
\textbf{comprehend complex ideas}, \textbf{learn quickly} and \textbf{learn from experience}.
It is not merely book-learning, a narrow academic skill, or test-taking smarts. Rather, it
reflects a broader and deeper capability for comprehending our surroundings:
``catching on'', ``making sense'' of things, or ``figuring out'' what to do.
\end{nbdef}

The definition was published in the \emph{Wall Street Journal} in December 1994 and signed
by 52 scientists. The professor remarks that the statement also contains claims she does not
feel comfortable with: it is a reference point, not an undisputed truth.
The seven abilities in bold (highlighted in red on slide~6) are worth remembering: each of
them has become a branch of AI.

\begin{center}
\small
\begin{tabular}{@{}ll@{}}
\toprule
\textbf{Ability in the definition} & \textbf{AI field that tries to reproduce it} \\
\midrule
reason                          & logic, automated reasoning, Prolog (this course) \\
plan                            & automated planning (e.g.\ the STRIPS planner of Shakey) \\
solve problems                  & search algorithms, constraint solving \\
think abstractly                & knowledge representation, ontologies \\
comprehend complex ideas        & natural language understanding \\
learn quickly / from experience & machine learning, reinforcement learning \\
\bottomrule
\end{tabular}
\captionof{table}{The abilities in the 1994 definition and the AI research areas that study them.}
\end{center}

\subsubsection{Many kinds of intelligence}
\slidesrc{Thinking machines: 7--8}

Intelligence is not a single number. Two influential theories:

\begin{nbdef}{Multiple intelligences (H.~Gardner, \emph{Frames of Mind}, 1983)}
Gardner proposed that humans have several distinct intelligences:
\begin{multicols}{2}
\begin{enumerate}[nosep]
  \item musical-rhythmic and harmonic;
  \item visual-spatial;
  \item verbal-linguistic;
  \item logical-mathematical;
  \item bodily-kinesthetic;
  \item interpersonal (understanding others);
  \item intrapersonal (understanding oneself);
  \item naturalistic (recognising and classifying nature);
  \item existential (asking deep questions about life).
\end{enumerate}
\end{multicols}
\end{nbdef}

\begin{nbdef}{Emotional intelligence (D.~Goleman, 1995)}
Emotional intelligence (EI) is the capability of individuals to recognise their own emotions
and those of others, discern between different feelings and label them appropriately, use
emotional information to guide thinking and behaviour, and manage and/or adjust emotions
to adapt to environments or achieve one's goal(s).
\end{nbdef}

\begin{nbintuition}[Why does an AI course care?]
If intelligence has many faces, then ``artificial intelligence'' can also mean many
different things. A chess program shows logical-mathematical intelligence; a robot that
walks on rough ground shows bodily-kinesthetic intelligence; a chatbot that adapts its tone
to an angry customer imitates a bit of emotional and interpersonal intelligence. A
\emph{multi-agent system} (the main topic of this course) needs \emph{interpersonal}
(social) intelligence: agents must understand and negotiate with each other.
\end{nbintuition}

\subsubsection{Thinking, fast and slow: System 1 and System 2}
\label{sec:system1-system2}
\slidesrc{Thinking machines: 9--11}

In the book \emph{Thinking, Fast and Slow} (2011), the psychologist and Nobel prize winner
Daniel Kahneman describes the mind as two cooperating ``systems''.

\begin{nbdef}{System 1 and System 2}
\begin{itemize}[nosep]
  \item \textbf{System 1}: fast, automatic, frequent, emotional, stereotypic, unconscious.
  \item \textbf{System 2}: slow, effortful, infrequent, logical, calculating, conscious.
\end{itemize}
\end{nbdef}

\begin{center}
\small
\begin{tabularx}{\linewidth}{@{}XX@{}}
\toprule
\textbf{System 1 can\ldots{}} (in order of complexity) & \textbf{System 2 is needed to\ldots{}} \\
\midrule
determine that an object is farther than another & prepare for the start of a sprint \\
localise the source of a sound & direct your attention to the clowns at the circus \\
complete the phrase ``war and \ldots'' & look for the woman with grey hair \\
show disgust at a gruesome image & try to recognise a sound \\
solve $2+2=?$ & sustain a faster-than-normal walking rate \\
read text on a billboard & judge if a behaviour is appropriate in a social setting \\
drive a car on an empty road & count the number of A's in a text \\
think of a good chess move (if you are a chess master) & give someone your phone number \\
understand simple sentences & park into a tight parking space \\
 & compare the price/quality ratio of two washing machines \\
 & determine the validity of a complex logical reasoning \\
\bottomrule
\end{tabularx}
\captionof{table}{Kahneman's examples (slides 10--11). Notice that the same task can move
from System~2 to System~1 with practice: a chess master ``sees'' a good move.}
\end{center}

\begin{nbexample}{the bat and the ball}
\notinslides
A bat and a ball cost 1.10 euro in total. The bat costs 1 euro more than the ball. How much
does the ball cost? System~1 shouts ``10 cents!''. System~2, if you wake it up, checks:
if the ball costs $x$, the bat costs $x+1$, so
\begin{align*}
x + (x + 1) &= 1.10 \\
2x &= 0.10 \\
x &= 0.05 .
\end{align*}
The ball costs 5 cents (check: $0.05 + 1.05 = 1.10$). Fast intuition is cheap but can be
wrong; slow reasoning is expensive but verifiable. Keep this in mind: it is exactly the
difference between \emph{sub-symbolic} and \emph{symbolic} AI
(Section~\ref{sec:symbolic-vs-subsymbolic}).
\end{nbexample}

\subsubsection{Artificial intelligence and intelligent agents}
\label{sec:ai-definition}
\slidesrc{Thinking machines: 12}

\begin{nbdef}{Artificial intelligence}
Artificial intelligence (AI) is intelligence demonstrated by machines, unlike the natural
intelligence displayed by humans and animals. The leading AI textbook, \emph{Artificial
Intelligence: A Modern Approach} (S.~Russell and P.~Norvig, 4th edition), defines the field
as the study of \textbf{intelligent agents}: any device that \textbf{perceives} its
environment and \textbf{takes actions} that maximise its chance of successfully achieving
its \textbf{goals}.
\end{nbdef}

This definition is the bridge to the whole course: AI is not defined by \emph{how} the
device works inside (neural network, logic rules, \ldots) but by its behaviour as an
\textbf{agent}, a loop of perceiving and acting towards goals.

\begin{center}
\begin{tikzpicture}[>=Stealth, node distance=1.2cm,
  box/.style={draw, rounded corners, thick, minimum width=2.6cm, minimum height=1cm, align=center}]
  \node[box, draw=nbblue, fill=nbblue!8] (agent) {\textbf{Agent}\\[-1pt]\footnotesize goals};
  \node[box, draw=nbgreen, fill=nbgreen!8, right=3.2cm of agent] (env) {\textbf{Environment}};
  \draw[->, thick, nbgreen!70!black] (env.north west) to[bend right=25] node[above, font=\small]{percepts (sensors)} (agent.north east);
  \draw[->, thick, nbred] (agent.south east) to[bend right=25] node[below, font=\small]{actions (actuators)} (env.south west);
\end{tikzpicture}
\captionof{figure}{An intelligent agent: it perceives the environment and acts on it to reach its
goals. Part~III of these notes is entirely about this picture.}
\end{center}

\begin{nbwhy}
Everyday intelligent agents: a robot vacuum cleaner (perceives obstacles and dirt, acts by
moving and sucking, goal: a clean floor); a smart thermostat (perceives temperature and your
habits, acts on the boiler, goal: comfort at minimum cost); a trading bot (perceives prices,
acts by buying/selling, goal: profit); a navigation app (perceives traffic, acts by
suggesting a route, goal: get you there fast).
\end{nbwhy}
